\section{A bounded approximation procedure for the weights}
\label{sec:evaluation}

The free large coordinates in \eqref{eq:hj} cannot be enumerated in
polynomial time.  We instead partition their ranges into polynomially
many intervals per coordinate and sample the resulting product grid.
Scaled free coordinates, softened constraints, and grid walks appear
in the dense-table methods of Dyer, Kannan, and Mount~\cite{DKM1997};
Morris~\cite{Morris2002} further developed the associated geometric
rounding approach. The weighted bin construction below includes its
geometric estimate, population correction, and finite-bit
implementation in full.

\begin{samepage}
\begin{proposition}[Weight evaluation]
\label{prop:weight-oracle}
Fix a box-valid small-count profile $(x,y)$ with total $p$, an integer
$0\le j\le H$, and rational parameters
$0<\xi_{\mathrm{in}},\theta<1/4$.  There is a randomized procedure
returning a positive rational $\widehat h$ such that
\[
 \mathbb P\left[
  (1-\xi_{\mathrm{in}})h_j(x,y)\le\widehat h
       \le(1+\xi_{\mathrm{in}})h_j(x,y)
 \right]\ge1-\theta.
\]
On every execution its bit cost is bounded by a fixed polynomial in
$L,\epsilon^{-1},\xi_{\mathrm{in}}^{-1},\theta^{-1}$ and the encoding
lengths of the parameters.  If
$M_f=\prod_{s\ \mathrm{free}}(b_s+1)$, every returned value satisfies
\begin{equation}
 \frac{M_f2^{-j}}{16}\le\widehat h\le8M_f.
 \label{eq:oracle-unconditional-range}
\end{equation}
In particular, positivity and the representation-length bound do not
depend on the statistical success event.
\end{proposition}
\end{samepage}

\subsection{Bins and a comparable continuous density}

Throughout the construction, fix $(x,y)$.  Let $e$ be the number of free
cells.  We first describe the positive-dimensional case; the case $e=0$
is addressed below.  On $[0,B]^e$, write $D(v)$ for the penalty of the
free values $z_s=(b_s+1)v_s/B$.  Lemma~\ref{lem:forest} gives
\begin{equation}
 0\le D(v)\le H,
 \qquad
 |D(v)-D(w)|\le\frac{2Ad^2}{B}\|v-w\|_\infty
                   =\frac18\|v-w\|_\infty.
 \label{eq:scaled-penalty-bounds}
\end{equation}
The function $D$ is convex.

Let $\mathcal V=\{0,\ldots,B-1\}^e$ be the bin indices.  Assign an
integer free vector $z$ to the bin whose coordinates are
$\lfloor Bz_s/(b_s+1)\rfloor$.  Denote the resulting Cartesian product
of integer intervals by $I_v$.  In a coordinate with $N=b_s+1$, the
interval for bin index $v_s$ is
\[
 \left\{\left\lceil\frac{Nv_s}{B}\right\rceil,
 \ldots,
 \left\lceil\frac{N(v_s+1)}{B}\right\rceil-1\right\}.
\]
Its cardinality differs from $N/B$ by at most one.  Since a free cell
is large, $N\ge W=100dB$, so every interval is nonempty.  Define
\[
 \alpha_v=\frac{B^e|I_v|}{M_f}.
\]
Each coordinate contributes a factor in
$[1-1/(100d),1+1/(100d)]$.  There are at most $d$ coordinates, and hence
\begin{equation}
 \frac12\le\alpha_v\le2.
 \label{eq:bin-populations}
\end{equation}
The intervals and $\alpha_v$ are computed by integer arithmetic.

For an integer $0\le a\le j$, put
\[
 V_a(v)=2^{-\lfloor aD(v)/H\rfloor},
 \qquad
 \varrho_a(w)=2^{-aD(w)/H}\quad(w\in[0,B]^e).
\]
The density $\varrho_a$ is positive and log-concave.  Its base-two
logarithm has sup-norm Lipschitz constant at most $1/8$, and therefore
also at most one.  Combining this weaker bound with the floor gives,
on the unit bin with lower corner $v$,
\begin{equation}
 \frac{V_a(v)}4\le\varrho_a(w)\le2V_a(v).
 \label{eq:bin-density-comparison}
\end{equation}
The same inequalities hold at boundary points by continuity.

\subsection{The estimator and its population correction}

For $0\le a\le j$, define the bin partition function and law
\[
 S_a=\sum_{v\in\mathcal V}V_a(v),\qquad
 \pi_a(v)=\frac{V_a(v)}{S_a}.
\]
The dyadic weights $V_a$ will permit exact finite-bit Metropolis
acceptance. Their normalizing constants can be estimated through
adjacent ratios. For $0\le a<j$, the observable
\[
 R_a(v)=\frac{V_{a+1}(v)}{V_a(v)}
\]
has expectation $S_{a+1}/S_a$ under $\pi_a$.  Its values are $1/2$ or
one: $0\le D/H\le1$ makes the increment of the floor exponent zero
or one.

For the final correction, draw $v$ from $\pi_j$, then a uniform integer
point $z\in I_v$, and let $s$ be its cube coordinates
$s_u=Bz_u/(b_u+1)$.  Use
\begin{equation}
 T(v,z)=\alpha_v\frac{2^{-jD(s)/H}}{V_j(v)}.
 \label{eq:oracle-correction}
\end{equation}
Equations~\eqref{eq:bin-populations} and
\eqref{eq:bin-density-comparison} show $1/8\le T\le4$.
The factor $\alpha_v$ corrects for the number of integer points in
the selected bin. The correction therefore has exact expectation
\[
 \mathbb E T
 =\frac{1}{S_j}\sum_v\frac{\alpha_v}{|I_v|}
                         \sum_{z\in I_v}2^{-jD(s)/H}
 =\frac{B^eh_j(x,y)}{M_fS_j}.
\]
As $S_0=B^e$, the identity
\begin{equation}
 M_f\left(\prod_{a=0}^{j-1}\mathbb E_{\pi_a}R_a\right)
       \mathbb E T=h_j(x,y)
 \label{eq:oracle-telescope}
\end{equation}
specifies the estimator.  We approximate each expectation by an average
of independent fresh samples and multiply the $j+1$ averages by $M_f$.

\subsection{Interpolation and conductance}

It remains to sample the bin laws $\pi_a$ with controlled error and
bounded work. We use a local Metropolis walk. The continuous density
$\varrho_a$ supplies the cut expansion estimate that bounds its mixing
time; it is not itself sampled. We first establish the finite-box form
of the Pr\'ekopa--Leindler interpolation
inequality~\cite{Leindler1972,Prekopa1973} used in the geometric
argument. All integrals in the following lemma are with
respect to Lebesgue measure.

\begin{lemma}[Integral interpolation]
\label{lem:integral-interpolation}
Let $S_0,S_1,S_2$ be finite unions of bounded open axis-aligned boxes
contained in a common compact box.  Let $0<t<1$ and suppose
$(1-t)S_1+tS_2\subseteq S_0$.  Let $f_0,f_1,f_2$ be positive continuous
functions on the ambient compact box such that
\[
 f_0((1-t)x+ty)\ge f_1(x)^{1-t}f_2(y)^t
 \qquad(x\in S_1,\ y\in S_2).
\]
Then
\[
 \int_{S_0}f_0\ge
 \left(\int_{S_1}f_1\right)^{1-t}
 \left(\int_{S_2}f_2\right)^t.
\]
\end{lemma}

\begin{proof}
In dimension one, write $A_i=\int_{S_i}f_i$ for $i=1,2$; zero source
mass makes the conclusion immediate.  Let $x(u),y(u)$ be the increasing
quantiles of the two normalized source densities, $0<u<1$.  Away from
finitely many subdivision parameters they are continuously
differentiable, with
\[
 x'(u)=\frac{A_1}{f_1(x(u))},
 \qquad y'(u)=\frac{A_2}{f_2(y(u))}.
\]
The map $z(u)=(1-t)x(u)+ty(u)$ is increasing.  Its smooth pieces have
disjoint images contained in $S_0$; jumps merely omit nonnegative
contributions to the target integral.  Change of variables on these
pieces, followed by the weighted arithmetic--geometric mean inequality,
gives
\begin{align*}
 \int_{S_0}f_0
 &\ge\int_0^1 f_0(z(u))
 \left((1-t)\frac{A_1}{f_1(x(u))}
             +t\frac{A_2}{f_2(y(u))}\right)\,du\\
 &\ge\int_0^1 A_1^{1-t}A_2^t\,du
  =A_1^{1-t}A_2^t.
\end{align*}
This argument also permits nonnegative piecewise-continuous densities
with finitely many zero intervals, provided each nonzero piece is
bounded above and bounded away from zero.

For the induction on dimension, fix first coordinates $x_1,y_1$ of
the two sources and put $z_1=(1-t)x_1+ty_1$.  The corresponding slices
satisfy the same set inclusion and interpolation hypothesis in one
lower dimension.  The induction hypothesis therefore gives the
one-dimensional interpolation inequality for the three marginal
integrals over slices.  Subdivide the first-coordinate axis at the
finitely many box endpoints.  On each resulting interval every slice
domain is constant; positivity and continuity on the ambient compact
box make each nonzero marginal continuous, bounded, and bounded away
from zero there.  Thus the one-dimensional argument just proved applies
to the marginals.  Iterated integration proves the claim.
\end{proof}

On $\mathcal V$, use the Metropolis chain with weights $V_a$.  Each
signed coordinate direction is proposed with probability
$2^{-\lceil\log_2(4d)\rceil}$.  Hold for unused proposal outcomes or
for proposals outside $\mathcal V$, and accept a valid proposal with
probability $\min\{1,V_a(w)/V_a(v)\}$.  The chain is reversible for
$\pi_a$ and is lazy, since the total probability
of directional proposals is at most $1/2$.

\begin{lemma}[Bin conductance]
\label{lem:bin-conductance}
For $e\ge1$, every nontrivial cut of this chain has stationary flow
divided by its smaller stationary side mass at least
\[
 \kappa=\frac{1}{10^4d^2B}.
\]
\end{lemma}

\begin{proof}
Write $\mu$ for unnormalized integration against $\varrho_a$ on
$D_0=(0,B)^e$.  For a fixed cut, let $S$ be the union of the open unit
bins on the side of smaller continuous mass.  Grid hyperplanes have
zero mass, so $\mu(S)\le\mu(D_0)/2$.  Set
\[
 t=\frac{1}{4dB},\qquad r=tB=\frac{1}{4d},
 \qquad M=(1-t)S+tD_0.
\]
The set $M$ is a finite union of open boxes and contains $S$.
Lemma~\ref{lem:integral-interpolation}, applied with all three
functions equal to $\varrho_a$, gives
\begin{equation}
 \mu(M\setminus S)\ge(2^t-1)\mu(S)\ge\frac t2\mu(S).
 \label{eq:bin-expansion}
\end{equation}

We cover the added mass by neighborhoods of cut faces.  Consider a
point $z\in M\setminus S$ off the grid hyperplanes, and choose
$x\in S$, $y\in D_0$ with $z=(1-t)x+ty$.  The segment from $x$ to $z$
has sup-norm length less than $r$.  At some grid-crossing event its
entering and exiting bins have different cut membership.  If several
coordinates cross simultaneously, flip those coordinates one at a time
between the entering and exiting bins.  All intermediate bins are
incident to the crossing point and lie in the cube.  At least one
consecutive pair belongs to opposite sides of the cut.  Its closed
common face contains the crossing point, so $z$ is within sup distance
$r$ of an actual axis-adjacent cut face.

This proves containment, up to the null set of grid hyperplanes, in the
union of the radius-$r$ neighborhoods of the cut faces.  It does not
charge a multiplicity for the number of incident bins.  Disconnectedness
of $S$ causes no change in the argument.  The segment stays in the open
ambient cube, and intersecting a face neighborhood with that cube only
reduces its volume.

Each face neighborhood has volume at most
\[
 2r(1+2r)^{e-1}
 =\frac{1}{2d}\left(1+\frac{1}{2d}\right)^{e-1}\le1.
\]
Every point in it is within sup distance $1+r$ of either adjacent
lower corner.  The logarithmic Lipschitz bound and
$\varrho_a(v)\le V_a(v)$ therefore bound its density by
$2^{1+r}V_a(v)\le4V_a(v)$ for each of the two adjacent corners
separately.  Thus, if $F$ is the sum of
$\min\{V_a(v),V_a(w)\}$ over cut edges, the covering gives
$\mu(M\setminus S)\le4F$.

Let $m$ be the smaller of the two unnormalized discrete side masses.
By \eqref{eq:bin-density-comparison},
\[
 \mu(S)\ge\frac14\sum_{v:\,v+(0,1)^e\subset S}V_a(v)\ge\frac m4.
\]
This remains true if the smaller continuous side is the larger discrete
side.  Combining with \eqref{eq:bin-expansion} yields $F\ge tm/32$.
The directional proposal probability is at least $1/(8d)$, so the
normalized conductance is at least
\[
 \frac{t}{256d}=\frac{1}{1024d^2B}\ge\kappa.
\]
\end{proof}

For completeness, this cut bound supplies a quantitative mixing time
without an additional sampling theorem.  For a finite reversible kernel
$P$ with stationary law $\pi$, write
\[
 \mathcal E_P(g)=\frac12\sum_{v,w}\pi(v)P(v,w)(g(v)-g(w))^2.
\]
For nonnegative $g$ supported on stationary mass at most $1/2$,
integration over superlevel sets of $g^2$ gives
\[
 \kappa\mathbb E_\pi g^2
 \le\sum_{\{v,w\}}\pi(v)P(v,w)|g(v)^2-g(w)^2|
 \le\bigl(2\mathcal E_P(g)\mathbb E_\pi g^2\bigr)^{1/2}.
\]
Here the sum is over unordered distinct pairs, and the last inequality
uses Cauchy--Schwarz and
$(g(v)+g(w))^2\le2g(v)^2+2g(w)^2$.
Apply the resulting inequality to the positive and negative parts about
a median of an arbitrary function.  Their energies sum to at most its
energy, and their squared norms sum to at least its variance.  Hence
the inverse spectral gap is at most
\[
 K_{\mathrm{bin}}=\frac{2}{\kappa^2}
                  =2(10^4d^2B)^2.
\]
Laziness puts the spectrum in $[0,1]$: $2P-I$ is a reversible stochastic
contraction.  Since $2^{-H}\le V_a\le1$,
$\pi_{a,\min}^{-1}\le2^HB^e$.  Spectral contraction of the centered
density of an initial point, followed by Cauchy--Schwarz, therefore
gives the following convenient loose bound after $\ell$ steps:
\begin{equation}
 \|P_a^\ell(v,\cdot)-\pi_a\|_{\mathrm{TV}}
 \le 2^{H+dB}\exp(-\ell/K_{\mathrm{bin}}).
 \label{eq:bin-mixing}
\end{equation}

\subsection{Sample counts and error allocations}

We now implement the estimator in \eqref{eq:oracle-telescope}.
Each occurrence of a bin sample uses a fresh walk; each offset sample
uses fresh bits conditional on its bin.

Set
\begin{equation}
 \sigma=\frac{\xi_{\mathrm{in}}}{50(H+1)},\qquad
 N_{\mathrm{in}}=
 \left\lceil\frac{10^6(H+1)}{\theta\sigma^2}\right\rceil,
 \qquad Q_{\mathrm{in}}=(d+1)(H+1)N_{\mathrm{in}}.
 \label{eq:inner-parameters}
\end{equation}
Use $N_{\mathrm{in}}$ samples for every average.  Each exact observable
has variance divided by squared mean at most $1024$.  Chebyshev's
inequality and a union bound over at most $H+1$ averages imply that
all ideal averages have relative error at most $\sigma$, except with
probability at most
\[
 \frac{1024(H+1)}{N_{\mathrm{in}}\sigma^2}<\frac\theta4.
\]

Each bin sample is obtained by an independent walk, started at a fixed
bin, of length
\begin{equation}
 \ell_{\mathrm{in}}=
 \left\lceil K_{\mathrm{bin}}
  \left(H+dB+\frac{16Q_{\mathrm{in}}}{\theta}+2\right)\right\rceil.
 \label{eq:inner-walk-length}
\end{equation}
Equation~\eqref{eq:bin-mixing} bounds its total-variation error by
$\theta/(8Q_{\mathrm{in}})$.  This deliberately conservative length
is polynomial in the allowed inverse failure parameter.

\begin{lemma}[Finite-bit uniform integer draw]
\label{lem:dyadic-integer-draw}
For a positive integer $M$ and $0<\zeta<1$, choose the smallest
nonnegative integer $k$ with $M2^{-k}\le\zeta$.  If $Y$ is uniform on
$\{0,\ldots,2^k-1\}$, then
$\lfloor MY/2^k\rfloor$ has total-variation distance at most $\zeta$
from the uniform law on $\{0,\ldots,M-1\}$.  The draw uses exactly
$k$ fair bits and bounded integer arithmetic.
\end{lemma}

\begin{proof}
Each output has either $\lfloor2^k/M\rfloor$ or
$\lceil2^k/M\rceil$ preimages.  Summing its deviation from probability
$1/M$ gives a total-variation bound at most $M2^{-k}$.  Computing the
displayed integer quotient needs no rejection step.
\end{proof}

Apply Lemma~\ref{lem:dyadic-integer-draw} separately to every scalar
offset interval, with $\zeta=\theta/(8Q_{\mathrm{in}})$, and translate
by its lower endpoint.  There are at most $Q_{\mathrm{in}}$ bin draws
and scalar offset draws together.  Couple each bin draw to a stationary
one, and, conditional on a matching bin, couple the offsets to exact
uniform offsets in that bin.  On a bin mismatch retain the correct
conditional offset marginals on both sides; its failure is already
charged.  Use independent fresh randomness for different samples.
The resulting ideal pairs $(v,z)$ are independent across samples with
joint law $\pi_j(v)\operatorname{Unif}(I_v)$, although bin and offset
within a pair need not be independent.  The probability of any draw
discrepancy is at most $\theta/8$, in particular less than $\theta/4$.

\subsection{Correction arithmetic and a bound on every execution}

The floors defining $V_a$ are evaluated by exact rational arithmetic.
All these weights are dyadic.  Bin proposals can be made with
$\lceil\log_2(4d)\rceil$ bits, with unused outcomes interpreted as
holds.  Each Metropolis acceptance is an exact dyadic trial with at
most $H$ additional bits.  Thus the walk length
\eqref{eq:inner-walk-length} is also a deterministic bound on the
number of transitions; no random-time rejection procedure is used.

Only the power of two in \eqref{eq:oracle-correction} needs numerical
approximation.  Its exponent is the rational number
\[
 \beta=\left\lfloor\frac{jD(v)}H\right\rfloor
                          -\frac{jD(s)}H\in[-2,2].
\]
Indeed the floor loses less than one and the exponent changes by at
most one inside a bin.  Choose a power-of-two integer $s'$ with
$32/\sigma\le s'\le64/\sigma$, and round $\beta$ to its nearest
grid point $h/s'$.  Since $-2$ and $2$ are grid points, $|h|\le2s'$,
and $|\beta-h/s'|\le\sigma/64$.  If $T_0=2^\beta$, then
\[
 \left|\frac{2^{h/s'}}{T_0}-1\right|
 \le e^{\sigma/64}-1\le\frac\sigma{32},
 \qquad T_0\ge\frac14.
\]
Here $\sigma<1/200$.  Bisect the interval $[0,8]$, comparing the
$s'$-th power of a rational midpoint with $2^h$, for a fixed number of
iterations sufficient to make the absolute error at most $\sigma/64$.
The returned rational $R$ then satisfies
\[
 \left|\frac R{T_0}-1\right|
 \le\frac\sigma{32}+\frac\sigma{16}<\sigma,
 \qquad R\ge\frac14-\frac\sigma{64}>0.
\]
Here the rounded exponent $h/s'$ still belongs to $[-2,2]$, so its
exact power is at least $1/4$ before the bisection error is applied.
Negative $h$ causes no difficulty: $2^h$ is an exact dyadic rational
with $O(|h|)$ bits.  For example, a depth of
$\lceil\log_2(512/\sigma)\rceil$ and the final interval midpoint
suffice.  The bisection depth is $O(\log(1/\sigma))$, and
the exact powers being compared have
$O(s'\log(1/\sigma))$ bits.  Since $s'=O(\sigma^{-1})$, all these
operations have polynomial bit cost.  Multiply $R$ by the exact
$\alpha_v$ to approximate the correction with deterministic pointwise
relative error at most $\sigma$.

On the coupled event in which every ideal empirical mean is accurate,
there are $j+1$ empirical multiplicative errors of at most $\sigma$
and one further such error from rounding the correction.  The bound
$|\log(1\pm\sigma)|\le2\sigma$ gives total absolute logarithmic error
at most
\[
 2\sigma(j+2)\le4\sigma(H+1)=\frac{2\xi_{\mathrm{in}}}{25}.
\]
For $0<\xi_{\mathrm{in}}<1/4$, this implies the relative interval
$[1-\xi_{\mathrm{in}},1+\xi_{\mathrm{in}}]$ in
Proposition~\ref{prop:weight-oracle}.  The exceptional probability is
less than $\theta/2$, and hence at most $\theta$.

These statistical estimates are not needed for termination or output
size.  On every execution, the ratio averages lie in $[1/2,1]$.
The exact correction lies in $[1/8,4]$, and its computed value lies in
$[(1-\sigma)/8,4(1+\sigma)]\subset[1/16,8]$.  Its average has the same
bounds.  This proves \eqref{eq:oracle-unconditional-range} on every
path, including paths where all empirical means are inaccurate.

To check bit complexity, every loop is bounded by
\eqref{eq:inner-parameters} and \eqref{eq:inner-walk-length}.
The parameters $A,B,W,H$ are polynomial in $L$ and $\epsilon^{-1}$,
and these loop bounds are fixed polynomials in those parameters,
$\xi_{\mathrm{in}}^{-1}$, and $\theta^{-1}$.  Bin indices have
$O(\log B)$ bits, and rational cube coordinates used in the correction
have $O(L+\log B)$ bits.  Integer offsets have the binary lengths of the
input capacities; a scalar offset of interval length $M$ needs only
$O(\log M+\log(Q_{\mathrm{in}}/\theta))$ bits in
Lemma~\ref{lem:dyadic-integer-draw}.  The tree values are affine sums
of input margins and the current coordinates.  Divisions by tree
capacities in $D$ produce rational denominators whose total binary
length is polynomial, since only polynomially many input integers
occur.  The quantities $M_f$, $B^e$, $|I_v|$, and $\alpha_v$ likewise
have polynomial-length representations.  Finally, sums and products
of polynomially many rationals of polynomial length have polynomial
length even without cancellation.  This includes all empirical means
and the product forming $\widehat h$.  Thus no numerical large
capacity enters a loop bound.

If $e=0$, take $\mathcal V$ to consist of the empty vector and use the
empty-product conventions $M_f=B^e=\alpha_v=1$.  There is no offset
draw and no bin walk.  The same ratios and correction telescope to the
single summand $2^{-jD/H}$; the deterministic rounding and all stated
bounds remain valid.  If $j=0$, there are no ratio averages and
$V_0=1$.  The unrounded correction is $\alpha_v$, whose expectation
under uniform bins is exactly one.  Its numerical approximation has
the already allocated pointwise tolerance, so the procedure estimates
$h_0=M_f$ within the claimed error.  These observations complete the
proof of Proposition~\ref{prop:weight-oracle}.
